\documentclass[10pt,a4paper]{amsart}
\usepackage[margin=19mm]{geometry}
\usepackage{amsmath,amssymb,mathtools}
\usepackage[scr=rsfs,cal=euler]{mathalfa}
\usepackage{xfrac}
\usepackage[unicode,hidelinks,bookmarks=false]{hyperref}
\newcommand{\ie}{i.e.\ }
\newcommand{\cf}{cf.\ }
\newcommand{\resp}{respectively\ }
\newcommand{\Subsection}{\subsection}
\providecommand{\nobreakdash}{\nobreak\hbox{-}}
\setlength{\emergencystretch}{2em}
\multlinegap0pt
\usepackage{bm}
\usepackage{bbm}
\usepackage{dsfont}
\usepackage{xspace}
\usepackage{cancel}
\usepackage{mathtools}
\usepackage{amsthm}
\makeatletter
\@ifundefined{theorem}{\newtheorem{theorem}{Theorem}}{}
\@ifundefined{lemma}{\newtheorem{lemma}[theorem]{Lemma}}{}
\@ifundefined{proposition}{\newtheorem{proposition}[theorem]{Proposition}}{}
\makeatother
\DeclareMathOperator{\dl}{dl}
\DeclareMathOperator{\lk}{lk}
\newcommand{\Pf}[1]{\mathcal{PF}\left(#1\right)}
\newcommand{\Pm}[1]{\mathcal{PM}\left(#1\right)}
\newcommand{\tup}[1]{\left( #1 \right)}
\title[The path-missing and path-free complexes of a directed graph]{The path-missing and path-free complexes of a directed graph}
\author{Darij Grinberg, Lukas Katth\"an, Joel Brewster Lewis}
\date{Source: arXiv:2102.07894v4 (CC BY 4.0)}
\begin{document}
\maketitle

\noindent\emph{Source and licence.} The path-missing and path-free complexes of a directed graph. Version arXiv:2102.07894v4 (CC BY 4.0), distributed under the Creative Commons Attribution 4.0 International licence, as recorded in the arXiv licence metadata for this exact version and shown on the abstract page https://arxiv.org/abs/2102.07894v4. Full rights notice: licenses/arxiv-2102.07894-CC-BY-4.0.txt.\par\smallskip

\noindent\textit{Editorial scope.} Counted unit: the Proposition whose source label is \texttt{prop.strong-grapes.PfPm} in the article (source0.tex lines 2921--2924, source label \texttt{prop.strong-grapes.PfPm}), delivered together with the article's own complete original proof of it (source0.tex lines 2926--3064, one \texttt{proof} environment of 139 lines). No mathematics is written, repaired or re-derived: statement and proof are the article's own text line for line. Required context retained uncounted: lem:dumbarc (lines 386--389); lem:delcon (lines 672--678); lem:delcon\_pf (lines 723--729); lem:induction\_step\_1 (lines 833--838); lem:induction\_step\_2b (lines 880--886). Every definition, display and label of those lines is retained unchanged. Omitted from this package and not redistributed: 1--385, 390--671, 679--722, 730--832, 839--879, 887--2920, 2925--2925, 3065--3880. Nothing inside a retained excerpt is omitted or altered: every retained body is delivered exactly as the article's lines are written, so no omission receipt of any kind accompanies this package. Citations: the retained excerpts carry no citation command at all, so no bibliography entry of the article is transcribed and no reference is redistributed. Reference adaptation: none was needed and none was made; every reference inside the delivered unit resolves inside it, so no unresolved reference or citation is produced. Printed slips kept literal: no printed word, symbol, hypothesis or number of the counted statement or of its proof was altered. Layout and class: the article's own class is \texttt{amsart} with packages including amsmath, fontenc, inputenc, geometry, lmodern, bm, amssymb, amsfonts, paralist, comment, babel, csquotes, biblatex, tabls; the wrapper is the lane's pinned \texttt{amsart} editorial wrapper at 10pt on A4, which restates the theorem-like environments the retained text needs and adds the article's own macro definitions that the retained text needs (\texttt{\textbackslash Pf}, \texttt{\textbackslash Pm}, \texttt{\textbackslash dl}, \texttt{\textbackslash lk}, \texttt{\textbackslash tup}), with extra wrapper packages (bm, bbm, dsfont, xspace, cancel, mathtools). The delivered numbering is the unit's own. Assets: no retained span contains a figure environment, a graphics-inclusion command, a PDF-page inclusion, a table or a TikZ picture; no image file is copied into this package, the package declares no asset dependency and no image file is redistributed with it. Licence: version arXiv:2102.07894v4 of this article is distributed under the Creative Commons Attribution 4.0 International licence, as recorded in the arXiv licence metadata for this exact version and shown on the abstract page https://arxiv.org/abs/2102.07894v4. A full rights notice is delivered at licenses/arxiv-2102.07894-CC-BY-4.0.txt. Characters: every retained excerpt is pure ASCII, so the delivered engine, XeLaTeX with its default Latin Modern text font, reports no missing character.
\begin{lemma}\label{lem:dumbarc}
    Assume that $G$ has a useless arc $e$.
    Then, both $\Pf{G}$ and $\Pm{G}$ are cones with apex $e$.
\end{lemma}

\begin{lemma}\label{lem:delcon}
    Let $e \in E$. Then:
    \begin{enumerate}
        \item[(a)] We always have $\lk_{\Pm{G}}(e) = \Pm{G \backslash e}$.
        \item[(b)] We have $\dl_{\Pm{G}}(e) = \Pm{G / e}$ if the source of $e$ is $s$.
    \end{enumerate}
\end{lemma}

\begin{lemma}\label{lem:delcon_pf}
	Let $e \in E$. Then:
	\begin{enumerate}
		\item[(a)] We always have $\dl_{\Pf{G}}(e) = \Pf{G \backslash e}$.
		\item[(b)] We have $\lk_{\Pf{G}}(e) = \Pf{G / e}$ if the source of $e$ is $s$.
	\end{enumerate}
\end{lemma}

\begin{lemma}\label{lem:induction_step_1}
	Let $e \in E$ be an arc whose source is $s$.
    Let $s'$ be the target of $e$.
    Assume that $e$ is not the only arc with target $s'$.
    Then, $G / e$ has a useless arc.
\end{lemma}

\begin{lemma}\label{lem:induction_step_2b}
	Let $e \in E$ be an arc whose source is $s$.
    Assume that $e$ is not useless.
    Let $s'$ be the target of $e$.
	Assume further that $\# E > 1$, and that $e$ is the only arc with target $s'$.
    Then, $G \backslash e$ has a useless arc.
\end{lemma}

\begin{proposition}
\label{prop.strong-grapes.PfPm}
Both $\Pf{G}$ and $\Pm{G}$ are strong grapes.
\end{proposition}

\begin{proof}
We proceed by strong induction on $\# E$.
Thus, we assume (as induction hypothesis) that both $\Pf{H}$ and $\Pm{H}$ are strong grapes whenever $H$ is a graph with fewer arcs than $G$.
We must now show that both $\Pf{G}$ and $\Pm{G}$ are strong grapes.
This is obvious if $\# E \leq 1$, so we WLOG assume that $\# E > 1$.

The case when $s = t$ is easy\footnote{%
Indeed, in this case, it is clear that any subset of $E$
contains an $s-t$-path, and thus we have
$\Pf{G} = \varnothing$ and $\Pm{G} = 2^E$.
Thus, all that needs to be proved is that both $\tup{E, 2^E}$
and $\tup{E, \varnothing}$ are strong grapes.
But this is straightforward to do by induction on $\# E$
(taking any $a \in E$ in the induction step).}.
Thus, we WLOG assume that $s \neq t$.

The case when $G$ has no $s-t$-path is easy\footnote{%
Indeed, in this case, it can easily be checked that
$\Pf{G} = 2^E$ and $\Pm{G} = \varnothing$.
Thus, all that needs to be proved is that both $\tup{E, 2^E}$
and $\tup{E, \varnothing}$ are strong grapes.
But this was done in the previous footnote.}.
Thus, we WLOG assume that $G$ has at least one $s-t$-path.
This $s-t$-path has at least one arc (since $s \neq t$).
Thus, its first arc is well-defined.
This first arc must have source $s$ and is not useless (since it
belongs to an $s-t$-path).
Hence, there exists an arc $e$ with source $s$ that is not useless.
Consider this arc $e$.
Let $s'$ be its target.

The graph $G / e$ has fewer arcs than $G$.
Hence, by the induction hypothesis, both
$\Pf{G / e}$ and $\Pm{G / e}$ are strong grapes.
Likewise, both
$\Pf{G \backslash e}$ and $\Pm{G \backslash e}$ are strong grapes.

Lemma~\ref{lem:delcon_pf} (a) yields
that $\dl_{\Pf{G}}(e) = \Pf{G \backslash e}$,
whereas
Lemma~\ref{lem:delcon_pf} (b) yields
that $\lk_{\Pf{G}}(e) = \Pf{G / e}$.
Likewise, Lemma~\ref{lem:delcon} (a) yields
that $\lk_{\Pm{G}}(e) = \Pm{G \backslash e}$,
whereas
Lemma~\ref{lem:delcon} (b) yields
that $\dl_{\Pm{G}}(e) = \Pm{G / e}$.

Recall that $\Pf{G / e}$ and $\Pf{G \backslash e}$ are
strong grapes.
In other words, $\lk_{\Pf{G}}(e)$ and $\dl_{\Pf{G}}(e)$
are strong grapes
(since $\dl_{\Pf{G}}(e) = \Pf{G \backslash e}$ and
$\lk_{\Pf{G}}(e) = \Pf{G / e}$).
Thus, if we can show that at least one of
$\lk_{\Pf{G}}(e)$ and $\dl_{\Pf{G}}(e)$ is a cone,
then we can conclude that $\Pf{G}$ is a strong grape
(by the definition of a strong grape).

% Recall that $\Pm{G \backslash e}$ and $\Pm{G / e}$ are
% strong grapes.
% In other words, $\lk_{\Pm{G}}(e)$ and $\dl_{\Pm{G}}(e)$
% are strong grapes
% (since $\lk_{\Pm{G}}(e) = \Pm{G \backslash e}$ and
% $\dl_{\Pm{G}}(e) = \Pm{G / e}$).
% Thus,
Likewise, if we can show that at least one of
$\lk_{\Pm{G}}(e)$ and $\dl_{\Pm{G}}(e)$ is a cone,
then we can conclude that $\Pm{G}$ is a strong grape
%(by the definition of a strong grape)
(since $\lk_{\Pm{G}}(e) = \Pm{G \backslash e}$ and
$\dl_{\Pm{G}}(e) = \Pm{G / e}$ are strong grapes).

We are in one of the following two cases:

\textit{Case 1:} The arc $e$ is not the only arc with target $s'$.

\textit{Case 2:} The arc $e$ is the only arc with target $s'$.

Let us first consider Case 1.
In this case, the arc $e$ is not the only arc with target $s'$.
Hence, Lemma~\ref{lem:induction_step_1} shows that
$G / e$ has a useless arc.
Let $e'$ be this useless arc.
Thus, Lemma~\ref{lem:dumbarc} (applied to $G / e$ and $e'$ instead
of $G$ and $e$) shows that
both $\Pf{G / e}$ and $\Pm{G / e}$ are cones with apex $e'$.
In other words,
both $\lk_{\Pf{G}}(e)$ and $\dl_{\Pm{G}}(e)$ are cones
(since $\lk_{\Pf{G}}(e) = \Pf{G / e}$
and $\dl_{\Pm{G}}(e) = \Pm{G / e}$).
%
In particular, at least one of
$\lk_{\Pf{G}}(e)$ and $\dl_{\Pf{G}}(e)$ is a cone
(namely, $\lk_{\Pf{G}}(e)$).
As we have seen above, this entails that $\Pf{G}$ is a strong grape.
%
Furthermore, at least one of
$\lk_{\Pm{G}}(e)$ and $\dl_{\Pm{G}}(e)$ is a cone
(namely, $\dl_{\Pm{G}}(e)$).
As we have seen above, this entails that $\Pm{G}$ is a strong grape.

We thus have shown that
both $\Pf{G}$ and $\Pm{G}$ are strong grapes.
This completes the induction step in Case 1.

Let us now consider Case 2.
In this case, the arc $e$ is the only arc with target $s'$.
%
Hence, Lemma~\ref{lem:induction_step_2b} yields that
$G \backslash e$ has a useless arc.
Let $e'$ be this useless arc.
Thus, Lemma~\ref{lem:dumbarc} (applied to $G \backslash e$ and $e'$ instead
of $G$ and $e$) shows that
both $\Pf{G \backslash e}$ and $\Pm{G \backslash e}$ are cones with apex $e'$.
In other words,
both $\dl_{\Pf{G}}(e)$ and $\lk_{\Pm{G}}(e)$ are cones
(since $\lk_{\Pm{G}}(e) = \Pm{G \backslash e}$
and $\dl_{\Pf{G}}(e) = \Pf{G \backslash e}$).
%
In particular, at least one of
$\lk_{\Pf{G}}(e)$ and $\dl_{\Pf{G}}(e)$ is a cone
(namely, $\dl_{\Pf{G}}(e)$).
As we have seen above, this entails that $\Pf{G}$ is a strong grape.
%
Furthermore, at least one of
$\lk_{\Pm{G}}(e)$ and $\dl_{\Pm{G}}(e)$ is a cone
(namely, $\lk_{\Pm{G}}(e)$).
As we have seen above, this entails that $\Pm{G}$ is a strong grape.

We thus have shown that
both $\Pf{G}$ and $\Pm{G}$ are strong grapes.
This completes the induction step in Case 2.

We have now completed the induction step in each of
the two Cases 1 and 2.
Hence, the induction step is complete,
and Proposition~\ref{prop.strong-grapes.PfPm} is proved.
\end{proof}

\end{document}
